\chapter{Construction d'une table d'analyse fiable}
\label{chap:preparation}

Le notebook 03 repart des 24 CSV et construit une table métropolitaine par usager.
L'ordre des opérations est important~: les clés sont harmonisées, le périmètre est
sélectionné, les champs sont préparés et les lieux sont résumés avant toute jointure
finale. Les fichiers bruts ne sont pas modifiés.

\section{Lecture et conservation de la provenance}
Toutes les colonnes sont d'abord lues comme du texte, sans conversion automatique de
chaînes telles que \code{N/A} en valeurs manquantes. Une convention de lecture identique
évite de traiter différemment deux années avant même d'avoir défini les règles.

Chaque enregistrement reçoit son rang source et le chemin du CSV dont il provient.
Ce rang commence à 1 après l'en-tête~; il désigne un enregistrement CSV et pas forcément
une ligne physique si un champ contient un retour à la ligne. Le manifeste conserve
la taille, les dimensions et l'empreinte SHA-256 des 24 entrées. Les empreintes sont
revérifiées avant export pour détecter un changement pendant l'exécution.

L'identifiant \code{id_observation} combine l'année et le rang de l'usager, par exemple
\code{2024_usager_000011}. Il est créé pour toutes les années, y compris celles disposant
d'un identifiant officiel. Il reste stable pour une même version et un même ordre du
fichier source. Les usagers aux attributs identiques conservent des identifiants différents.

\section{Harmonisation et sélection géographique}
La clé \code{Accident_Id} de 2022 est renommée \code{Num_Acc}. Les espaces, y compris
internes et insécables, sont retirés des identifiants dans toutes les tables concernées.
Les clés et les collisions sont contrôlées avant de poursuivre. Le passage d'un code
de département comme 1 à 01 est explicite~; les codes 2A et 2B restent distincts.

La métropole est sélectionnée par une liste de 96 codes départementaux~: 01 à 95 hors
ancien code 20, auxquels s'ajoutent 2A et 2B. Le filtre repose sur le département déclaré
de l'accident, puis conserve les lieux, véhicules et usagers liés à ces accidents.
Un code inattendu provoquerait un arrêt pour examen, au lieu d'être silencieusement exclu.

Les codes commune sont aussi contrôlés. Un format non interprétable devient manquant
dans la colonne préparée, avec la valeur brute conservée. Une commune de format valide
dont le préfixe diffère du département est maintenue avec une alerte. Les quatre
accidents concernés par ces contrôles représentent six lignes d'usagers après jointure~:
le comptage d'une anomalie doit donc toujours préciser son unité.

Le tableau~\ref{tab:population-preparee} décrit le passage de 745\,158 usagers des sources
complètes à 702\,823 usagers métropolitains. Les 42\,335 lignes hors périmètre ne sont
pas des suppressions motivées par la gravité ou la qualité des variables.
\begin{table}[htbp]
\centering\small
\renewcommand{\arraystretch}{1.18}
\caption{Passage des sources complètes à la table métropolitaine.}\label{tab:population-preparee}
\begin{tabularx}{\textwidth}{@{}l>{\raggedleft\arraybackslash}X>{\raggedleft\arraybackslash}X>{\raggedleft\arraybackslash}X>{\raggedleft\arraybackslash}X@{}}
\toprule
Année & Usagers sources & Usagers retenus & Accidents retenus & Gravités inconnues retenues \\
\midrule
2019 & 132\,977 & 126\,352 & 56\,016 & 0 \\
2020 & 105\,295 & 99\,169 & 45\,121 & 0 \\
2021 & 129\,248 & 122\,189 & 53\,540 & 54 \\
2022 & 126\,662 & 119\,638 & 52\,380 & 215 \\
2023 & 125\,789 & 118\,197 & 51\,641 & 96 \\
2024 & 125\,187 & 117\,278 & 51\,058 & 0 \\
\textbf{Total} & 745\,158 & 702\,823 & 309\,756 & 365 \\
\bottomrule
\end{tabularx}
\sourcecalcul{notebook 03, bilan des effectifs}
\end{table}


\section{Règles pour les codes et les mesures}
Les colonnes catégorielles préparées contiennent un code documenté et interprétable,
ou une valeur manquante. Les variantes d'écriture telles que \code{1.0} sont ramenées
à \code{1}, mais le résultat reste du texte. Les lettres présentes dans certaines
rubriques sont préservées lorsque leur sens est documenté.

\begin{table}[htbp]\centering\small\renewcommand{\arraystretch}{1.2}
\caption{Exemples de règles appliquées aux valeurs particulières.}\label{tab:codes}
\begin{tabularx}{\textwidth}{@{}LLL@{}}\toprule
Situation & Interprétation & Traitement\\\midrule
\code{trajet = 0} & Non renseigné & Valeur préparée manquante.\\
\code{secu1 = 0} & Aucun équipement & Modalité conservée.\\
\code{vosp = 0} & Sans objet & Modalité conservée.\\
\code{actp = 0} & Non renseigné ou sans objet & Modalité mixte conservée, sans précision inventée.\\
Code non décrit, par exemple \code{actp = 7} & Sens non établi par la rubrique utilisée & Valeur préparée manquante, code source accessible.\\
\code{grav = -1} & Gravité inexploitable & Cible manquante, usager conservé.\\\bottomrule
\end{tabularx}
\sourcecalcul{règles du notebook 03, fondées sur la documentation et les valeurs observées}
\end{table}

Les codes \code{-1} sont laissés manquants~; lorsqu'ils ne figurent pas dans la liste
documentée d'un champ, ils sont aussi signalés comme non documentés dans le bilan.
Les colonnes brutes ou les CSV d'origine permettent de retrouver l'information initiale.
Les motifs de nettoyage sont comptés par variable et par table. Ils ne doivent pas être
additionnés pour produire un nombre de personnes~: une même ligne peut présenter
plusieurs difficultés et certaines tables décrivent des véhicules ou des lieux.

Le nombre de voies \code{nbv} et la vitesse maximale autorisée \code{vma} sont convertis
en mesures numériques. Une valeur vide, non numérique ou non positive devient manquante.
Un nombre de voies non entier est également inexploitable. Les autres valeurs positives
sont conservées. Les \code{vma} supérieures à 130 sont signalées pour examen, sans
correction supposée de 500 en 50. L'alerte porte sur 202 usagers associés à un accident
dont au moins un lieu dépasse ce seuil. Elle ne désigne ni 202 accidents ni une vitesse
réelle de véhicule~: \code{vma} est une limitation autorisée.

\section{Résumer les lieux sans inventer un lieu individuel}
Le résumé est construit séparément pour chaque variable au niveau de l'accident.
S'il reste une seule valeur connue distincte, elle est conservée. S'il existe plusieurs
valeurs connues, la valeur résumée est manquante et un indicateur signale le désaccord.
L'information partielle est un second indicateur, activé lorsqu'une valeur connue
coexiste avec une information inexploitable.

\begin{table}[htbp]\centering\small
\caption{Règle de résumé illustrée avec une limitation de vitesse.}\label{tab:consensus}
\begin{tabularx}{\textwidth}{@{}LLLL@{}}\toprule
Valeurs préparées & Résumé & Désaccord & Partielle\\\midrule
50, 50 & 50 & Non & Non\\
50, manquante & 50 & Non & Oui\\
50, 70 & Manquante & Oui & Non\\
50, 70, manquante & Manquante & Oui & Oui\\
Toutes manquantes & Manquante & Non & Non\\\bottomrule
\end{tabularx}
\sourcecalcul{convention du notebook 03 ; les indicateurs sont distincts}
\end{table}

Les champs résumés portent le préfixe \code{lieu_}. La règle est indépendante de l'ordre
des lignes et ne change pas si une répétition exacte est retirée. Pour l'accident
\code{202400000005}, la limitation résumée devient manquante avec désaccord, et les quatre
usagers restent quatre. En métropole, \code{vma} présente un désaccord entre valeurs
connues pour 1\,675 accidents en 2023 et 1\,875 en 2024. Ces nombres diffèrent des
variations brutes de l'audit, qui comptaient aussi les oppositions entre connu et inconnu.

Ce résumé est un compromis conservateur, pas une reconstitution du trajet de chaque
usager. Les lieux sources sont exportés séparément. Le nombre de lignes et les indicateurs
de désaccord servent d'abord à l'audit~; leur forte évolution entre années devra être
examinée avant tout usage prédictif.

\section{Âge, dates, coordonnées et cible}
L'âge approximatif est la différence entre l'année de l'accident et l'année de naissance.
L'anniversaire étant inconnu, il peut différer d'un an de l'âge exact. Les naissances
absentes, non interprétables ou non entières restent distinctes des âges hors intervalle.
La convention retenue conserve un âge préparé de 0 à 110 ans. Au-delà, l'âge devient
manquant, mais l'usager, l'estimation brute et l'indicateur d'alerte sont conservés.
Ce seuil est un outil de contrôle, pas une borne biologique démontrée.

Les 191 âges suspects des sources complètes deviennent 172 en métropole, dont 162 en
2019 et 10 en 2020. Le tableau~\ref{tab:qualite-preparee} présente les âges manquants
après traitement. Aucun âge n'est remplacé ici par une moyenne.
\begin{table}[htbp]
\centering\small
\renewcommand{\arraystretch}{1.18}
\caption{Quelques indicateurs de qualité dans le périmètre métropolitain.}\label{tab:qualite-preparee}
\begin{tabularx}{\textwidth}{@{}lLLL@{}}
\toprule
Année & Âges suspects & Âges manquants après traitement & Véhicules sans usager \\
\midrule
2019 & 162 & 162 & 1\,972 \\
2020 & 10 & 10 & 1\,838 \\
2021 & 0 & 2\,914 & 6 \\
2022 & 0 & 2\,732 & 0 \\
2023 & 0 & 2\,446 & 28 \\
2024 & 0 & 2\,419 & 23 \\
\bottomrule
\end{tabularx}
\sourcecalcul{notebook 03 ; les trois colonnes ne portent pas toutes sur la même unité}
\end{table}


Les dates sont converties et doivent correspondre à l'année du dossier~; les heures
doivent être interprétables. Des variables de mois, de jour de semaine et d'heure sont
créées. Les coordonnées utilisent un séparateur décimal harmonisé et un contrôle des
limites mondiales. Ces vérifications ne valident pas leur exactitude géographique.

La cible vaut 1 pour \code{grav} dans \{2,3\}, 0 pour \{1,4\}, et reste manquante sinon.
Les quatre catégories préparées et la gravité source restent accessibles. Les 419 gravités
inconnues des sources deviennent 365 en métropole~; la différence de 54 provient du
périmètre géographique. Ces 365 usagers sont conservés dans la table, mais ne devront
pas participer à un ajustement ou à une évaluation supervisés.

\section{Jointures et contrôles de conservation}
Les usagers constituent la table de départ. Le véhicule est ajouté sur
\code{Num_Acc} + \code{id_vehicule}, puis les caractéristiques et le résumé des lieux sur
\code{Num_Acc}. Chaque table de droite doit être unique sur sa clé. La validation
\code{many_to_one} interdit une multiplication des observations.

Avant et après chaque jointure, le programme compare l'effectif et l'ensemble des
\code{id_observation}. Il contrôle aussi les correspondances manquantes et l'accord des
deux versions de \code{num_veh}. Les 18 jointures annuelles conservent exactement leurs
usagers, sans correspondance manquante ni contradiction de numéro de véhicule.

Les champs propres aux piétons ne seront pas interprétés comme des comportements
piétons chez les autres usagers. Les trois équipements \code{secu1}, \code{secu2} et
\code{secu3} restent séparés~: aucun indicateur unique de protection n'est fabriqué
sans examiner leur contexte.

\section{Résultat, exports et relecture}
La table finale comporte \textbf{702\,823 lignes et 119 colonnes}, pour 309\,756 accidents.
Elle contient 702\,458 cibles connues et 365 inconnues. Aucun usager du périmètre n'a
été supprimé pour son âge, sa cible inconnue ou une variable explicative manquante.
Les 119 colonnes comprennent des identifiants, des valeurs sources, de la provenance
et des indicateurs de qualité~: elles ne forment pas une liste de 119 prédicteurs validés.

Deux CSV compressés sont produits~: la table des usagers et les lieux sources. Des JSON
décrivent les règles, les types, les rôles des colonnes et les empreintes. Six CSV de
bilan détaillent les effectifs, la géographie, les identifiants, les jointures, le
nettoyage et les lieux. L'annexe~\ref{ann:reproduction} en donne l'organisation.

Le CSV ne conserve pas les types Python. Le marqueur réservé
\code{__BAAC_MANQUANT__} distingue une valeur préparée manquante d'une chaîne vide
conservée dans un champ brut. Son absence dans les valeurs sources est vérifiée.
Une relecture par blocs applique le schéma, le marqueur et les conversions de date,
puis compare les 702\,823 lignes, toutes les colonnes, leurs valeurs et leurs types.

\FloatBarrier
\bilan{La table est exploitable sans confondre absence d'information et classe négative.
Les contrôles garantissent la conservation et la traçabilité des observations dans
les fichiers examinés. Ils ne corrigent pas les biais de collecte ni toutes les erreurs
de renseignement possibles.}
